\documentclass[12pt,a4paper]{article}
\usepackage{estatuto-config}


%==== VARIÁVEIS DO DOCUMENTO ====
\newcommand{\NomeDoCA}{[NOME DO CENTRO ACADÊMICO]}
\newcommand{\SiglaDoCA}{[SIGLA DO C.A.]}
\newcommand{\DataFundacao}{[Data da Fundação]}
\newcommand{\TipoDoCurso}{[Graduação ou Pós-Graduação]}
\newcommand{\NomeDoCurso}{[Nome do Curso]}
\newcommand{\NomeDaInstituicao}{[Nome da Universidade/Faculdade]}
\newcommand{\Cidade}{[Cidade]}
\newcommand{\Estado}{[Estado]}
\newcommand{\Instituto}{[Instituto de Ciências Exatas / Escola de Engenharia]}
%================================


\begin{document}

\begin{center}
    \Large\bfseries ESTATUTO DO CENTRO ACADÊMICO\\[0.2em]
    \normalsize\itshape \NomeDoCA{} DA UNIVERSIDADE FEDERAL DE MINAS GERAIS
\end{center}

\begin{center}
    \textit{"You can do imperfect work in an imperfect system and still make a difference."}
\end{center}

\vspace{1em}

\todo[inline]{\textbf{Comentário:} Caixas de texto com fundo laranja (como esta) não são parte do documento. São apenas para contextualizar ou comentar.}
\todo[inline,backgroundcolor=blue!70]{\textbf{Referência:} Caixas de texto com fundo azul (como esta) indicam uma referência a ser ajustada, uma vez o texto esteja com a composição de artigos final.}
\todo[inline,backgroundcolor=purple!30]{\textbf{Update:} Caixas de texto com fundo roxo claro (como esta) indicam onde houve alteração da última versão.}

%===================== CAPÍTULO I =====================
\capitulo{DA DENOMINAÇÃO, SEDE, DURAÇÃO E FINALIDADES}

\artigo{
O Centro Acadêmico \textit{\NomeDoCA{}}, doravante designado pela sigla \textit{\SiglaDoCA{}}, fundado em \textit{\DataFundacao{}}, é a entidade representativa do Corpo Discente do Curso de \textit{\TipoDoCurso{}} em \textit{\NomeDoCurso{}} da \textit{\NomeDaInstituicao{}}.

\paragrafoUnico{Tem sede e foro na cidade de \textit{\Cidade{}}, \textit{\Estado{}}, no endereço Avenida Antônio Carlos, 6627, \textit{\Instituto{}}, UFMG, Pampulha, Belo Horizonte – MG.}
}

\todo[inline]{\textbf{Comentário:} Aqui havia: \textbf{Art. 2} - O \textit{\SiglaDoCA{}} reconhece o Conselho Discente de Computação como instância de articulação entre os Centros Acadêmicos dos cursos da área de Computação. \textit{O referido conselho é composto, de forma nata, pelo Presidente ou representante oficialmente designado de cada centro acadêmico dos cursos diretamente ofertados
pelo [DCC] da UFMG.}}
\todo[inline]{\textbf{Comentário:} Sobre isso de reconhecer, parece haver uma tradição, sem nenhum valor aparente, de os estatutos de organizações na UFMG “reconhecerem” outras organizações. Isto não foi encontrado em estatutos de outras universidades públicas, em pesquisa rápida.}
\todo[inline]{\textbf{Comentário:} Se este trecho apenas “reconhece” algo de forma romântica, sem gerar efeitos práticos (jurídicos, orçamentários ou representativos), ele aumenta o Estatuto sem necessidade. Sugere-se retirar. A articulação pode seguir ocorrendo na prática.}

\artigo{
O \textit{\SiglaDoCA{}} é uma organização sem fins lucrativos, de duração indeterminada, laica e autônoma em relação à reitoria, à coordenação do curso, a partidos políticos e a outras entidades.
}

\artigo{
São finalidades do \textit{\SiglaDoCA{}}:
\inciso{
\item Representar o Corpo Discente do curso de \textit{\NomeDoCurso{}}, sem qualquer distinção de origem, raça, cor, etnia, nacionalidade, idade, condição de deficiência, condição socioeconômica, religião, sexo, identidade gênero e orientação sexual;
\todo[inline,backgroundcolor=purple!30]{\textbf{Update:} Adicionados origem, etnia, idade, condição socioeconômica. Alterado "gênero" para "identidade de gênero."}
\item Colaborar no aprimoramento do ensino, da pesquisa e da extensão universitária;
\item Incentivar a integração, o congraçamento e a solidariedade entre os associados;
\item Colaborar, de acordo com as possibilidades, com eventos de natureza acadêmica, científica, cultural, social e desportiva;
}
}

%===================== CAPÍTULO II =====================
\capitulo{DOS ASSOCIADOS, SEUS DIREITOS E DEVERES}

\artigo{
São associados natos do \textit{\SiglaDoCA{}} todos os estudantes regularmente matriculados no curso de \textit{\TipoDoCurso{}} em \textit{\NomeDoCurso{}} da UFMG.
}

\artigo{
São direitos dos associados:
\inciso{
\item Votar para os cargos da Diretoria Executiva, respeitadas as normas deste Estatuto e do respectivo edital de convocação;
\item Ser votado para os cargos da Diretoria Executiva, respeitadas as normas deste Estatuto e do respectivo edital de convocação;
\todo[inline,backgroundcolor=purple!30]{\textbf{Update:} O termo "regimento eleitoral" foi substituído por edital de convocação, por consistência, dado que era mencionado apenas aqui.}
\item Participar das Assembleias Gerais com direito a voz e a voto para definição de encaminhamentos e elaboração de propostas;
\item Apresentar propostas, sugestões e críticas à Diretoria e às Assembleias;
\item Usufruir dos serviços e benefícios oferecidos pelo \textit{\SiglaDoCA{}};
\item Ter acesso aos editais, às atas de reuniões, de assembleias e de plebiscitos, aos relatórios das comissões eleitorais e às prestações de contas financeiras, abrangendo também os de gestões anteriores;
\todo[inline,backgroundcolor=purple!30]{\textbf{Update:} Antes era "Ter acesso às atas, balancetes financeiros e demais documentos da gestão;". O termo "demais documentos" era vago.}
\item Votar e ser votado para as representações discentes em órgãos colegiados da UFMG, quando a escolha couber aos discentes do curso, conforme este Estatuto e as normas da UFMG;
}
}

\artigo{
São deveres dos associados:
\inciso{
\item Cumprir e fazer cumprir as disposições deste Estatuto;
\item Acatar as decisões aprovadas em Plebiscito;
\item Zelar pela preservação do patrimônio do \textit{\SiglaDoCA{}};
\item Indenizar todo e qualquer prejuízo ocasionado ao \textit{\SiglaDoCA{}} quando praticado intencionalmente;
}
}

\artigo{
A perda da condição de associado nato ocorrerá:
\inciso{
\item Automaticamente, pela perda do vínculo de matrícula com o curso de \textit{\TipoDoCurso{}} em \textit{\NomeDoCurso{}};
\item Por exclusão disciplinar, processada nos termos do Art. [8], aplicada estritamente no caso em que o associado impedir intencionalmente que outros associados exerçam os direitos garantidos no Art. [5] deste Estatuto.
}
\todo[inline,backgroundcolor=blue!70]{\textbf{Referência:} Art. [8] - O processo de exclusão disciplinar de associados garantirá o direito à ampla defesa e ao recurso [...]}
\todo[inline,backgroundcolor=blue!70]{\textbf{Referência:} Art. [5] - São direitos dos associados: [...]}
\todo[inline,backgroundcolor=purple!30]{\textbf{Update:} Antes era "Por exclusão disciplinar, em caso de infração grave ao Estatuto ou dano deliberado ao patrimônio."}
\todo[inline]{\textbf{Comentário:} A exclusão disciplinar foi restrita. O CA é formado por estudantes sem treinamento jurídico ou competência institucional para conduzir inquéritos. Transformar o CA em um "tribunal estudantil" abre precedentes para injustiças governadas por vieses e calúnias irresponsáveis. Crimes e infrações graves devem ser julgados pela UFMG ou pela Justiça (que já causam exclusão automática caso a matrícula seja cassada). A exclusão interna serve apenas para proteger a operação democrática do próprio CA.}
\paragrafoUnico{A exclusão disciplinar, mesmo se concretizada, não revoga os direitos aos serviços e benefícios básicos previstos no Art. 5º, inciso V.}
\todo[inline,backgroundcolor=blue!70]{\textbf{Referência:} Art. 5º, inciso V - Que dispõe sobre "Usufruir dos serviços e benefícios básicos".}
}

\artigo{
O processo de exclusão disciplinar de associados garantirá o direito à ampla defesa e ao recurso.
\inciso{
\item A pena de exclusão será aplicada, em primeira instância, por decisão formal e fundamentada da Diretoria Executiva, mediante notificação do associado afetado;
\item O associado poderá, mediante ato formal e expresso, apresentar sua defesa e recorrer desta decisão;
\item O recurso terá efeito suspensivo imediato, preservando os direitos do associado até a decisão final;
\item Havendo recurso, a Diretoria Executiva convocará, nos prazos regimentais, Assembleia Extraordinária para registrar a defesa, debater o caso e formular parecer a ser votado em Plebiscito.
}
\paragrafo{É vedada a aplicação de exclusão disciplinar contra membros da Diretoria Executiva.}
\todo[inline,backgroundcolor=purple!30]{\textbf{Update:} Adicionada proteção para impedir que a exclusão disciplinar seja usada para destituir diretores, burlando a exigência de plebiscito.}
\paragrafo{É vedada a divulgação de dados pessoais sem fundamento legítimo, cabendo à Diretoria decidir pela anonimização ou pela exigência de consentimento do associado.}
\todo[inline,backgroundcolor=purple!30]{\textbf{Update:} Antes era: "O associado poderá apresentar sua defesa e recorrer da decisão a qualquer momento;". A expressão "a qualquer momento" foi retirada para deixar o texto mais enxuto, pois a ausência de prazo já garante implicitamente que a pessoa recorra quando tomar ciência. A adição de "mediante ato formal e expresso" previne que conversas informais sejam consideradas como o exercício definitivo da defesa.}
\todo[inline]{\textbf{Comentário:} Como garantir a urbanidade, o respeito e a integridade das partes durante o debate na Assembleia? Prever moderação é delicado.}
}

%===================== CAPÍTULO III =====================
\capitulo{DA ESTRUTURA ORGANIZACIONAL}

\artigo{
São órgãos do \textit{\SiglaDoCA{}}:
\inciso{
\item A Diretoria Executiva.
\item A Assembleia Geral;
\item O Plebiscito;
}
}

%\artigo{É assegurada a participação remota dos associados nas Assembleias Gerais e nos Plebiscitos, bem como nos processos eleitorais, de forma simplificada e acessível, mediante autenticação necessária e suficiente para comprovar a identidade estudantil junto à universidade, vedada qualquer exigência cumulativa, excessiva ou sem relação direta com essa finalidade.\todo[inline]{Ponto sensível. Exigências documentais excessivas já ocorreram e desincentivam a participação. Ao mesmo tempo, há um desafio prático, pois a instituição não disponibiliza mais a lista de e-mails dos alunos.}\paragrafoUnico{Caberá aos responsáveis pela convocação da respectiva instância, ou à Comissão Eleitoral no caso de eleições, fornecer o suporte técnico necessário para garantir a acessibilidade dos participantes virtuais.}}

\artigo{
A oferta de participação e votação online é obrigatória em todas as Assembleias Gerais, Plebiscitos e eleições. A opção presencial é sempre opcional.
\paragrafo{O acesso deve ser simples e acessível, bastando uma autenticação que comprove a identidade estudantil junto à universidade. É proibida qualquer exigência acumulativa, excessiva ou que não cumpra diretamente com essa finalidade.}
\todo[inline]{\textbf{Comentário:} Ponto sensível. Exigências documentais excessivas já ocorreram e desincentivam a participação. Ao mesmo tempo, há um desafio prático, pois a instituição não disponibiliza mais a lista de e-mails dos alunos.}
\paragrafo{Cabe a quem convoca ou organiza fornecer o suporte necessário para assegurar a acessibilidade de todos à participação online.}
\todo[inline]{\textbf{Comentário:} A ideia era a de o voto ser possibilitado em computadores da biblioteca, por exemplo. Mas é positivo manter amplo para que se inove nas soluções conforme a necessidade.}
}

%===================== SEÇÃO I =====================

\secao{Da Diretoria Executiva}

\artigo{A Diretoria Executiva é o órgão responsável pela administração do \textit{\SiglaDoCA{}}. Seu mandato será de 1 (um) ano, permitida a reeleição de cada membro, por, no máximo, 2 (duas) vezes, independentemente do cargo ocupado, estendendo-se até a posse da nova chapa.
}

\artigo{A Diretoria Executiva será composta pelos seguintes cargos:
\inciso{
\item Presidência;
\item Secretaria;
\item Tesouraria;
\item Representação do Tronco Comum.
}
}

\artigo{São atribuições de cada cargo:
\inciso{
\item À Presidência: representar o \textit{\SiglaDoCA{}}, convocar e presidir reuniões, assinar documentos e movimentar contas juntamente com a Tesouraria;
\item À Tesouraria: gerenciar as finanças, elaborar os relatórios financeiros de prestação de contas e movimentar contas juntamente com a Presidência;
\todo[inline,backgroundcolor=purple!30]{\textbf{Update:} Updates: Antes era "elaborar balancetes e relatórios financeiros [...]".}
\item À Secretaria: redigir as atas das reuniões e de assembleias, bem como organizar e guardar os editais, as referidas atas, os relatórios financeiros e eleitorais e a correspondência oficial;
\todo[inline,backgroundcolor=purple!30]{\textbf{Update:} Antes era "redigir e guardar as atas, organizar os documentos e a correspondência oficial;".}
\item À Representação do Tronco Comum: Fazer serem conhecidas as demandas dos associados do Tronco Comum em Computação nas decisões da Diretoria Executiva.
\todo[inline]{\textbf{Comentário:} Os alunos do tronco comum ainda precisarão, de forma independente, votar e concordar em fazer parte do CA.}
}
}

\artigo{
Fica automaticamente desligado da Diretoria Executiva, perdendo o mandato, o associado que perder o vínculo de matrícula com o curso de \textit{\TipoDoCurso{}} em \textit{\NomeDoCurso{}} da UFMG.
}

\artigo{
A destituição de membros da Diretoria Executiva, seja ela parcial ou total, ocorrerá exclusivamente por meio de plebiscito, convocado a partir de proposta aprovada em Assembleia Geral Extraordinária.
\paragrafoUnico{O membro destituído ficará inelegível para a Diretoria Executiva nas 3 (três) eleições subsequentes.}
}

\todo[inline,backgroundcolor=purple!30]{\textbf{Update:} Os trechos abaixo foram atualizados com pequenas alterações.}
\artigo{
Em caso de vacância total da Diretoria Executiva, uma Assembleia Geral Extraordinária deverá ser convocada para, em caráter deliberativo excepcional e com dispensa de plebiscito, formar uma Comissão Eleitoral com a finalidade exclusiva de realizar um novo processo eleitoral emergencial.
\todo[inline,backgroundcolor=purple!30]{\textbf{Update:} A Assembleia agora possui esta competência excepcional (sem plebiscito) restrita à vacância total para evitar travamento institucional.}
\todo[inline,backgroundcolor=blue!70]{\textbf{Referência:} "(...) não possuindo caráter deliberativo, salvo nas exceções expressamente previstas neste Estatuto." / "Deliberar, em caráter excepcional e com dispensa de plebiscito, sobre a formação de Comissão Eleitoral em casos de vacância total da gestão, conforme os termos previstos neste Estatuto."}
\paragrafo{A vacância total será declarada caso ocorra qualquer uma das seguintes hipóteses:}
\inciso{
\item destituição de todos os membros da Diretoria;
\item renúncia coletiva da Diretoria sem a convocação prévia de eleições antecipadas, conforme facultado pelas normas do Processo Eleitoral;
\todo[inline,backgroundcolor=blue!70]{\textbf{Referência:} "A Diretoria Executiva poderá antecipar as eleições e o término do seu mandato.".}
\todo[inline,backgroundcolor=purple!30]{\textbf{Update:} A Diretoria pode encerrar o mandato antecipadamente publicando o edital de eleições. A vacância total da Diretoria como emergência só é acionada se eles renunciarem coletivamente sem sequer terem convocado as novas eleições.}
\item ausência de propostas de chapa no pleito regular;
\item decurso de 2 (dois) semestres letivos inteiros sem Diretoria ou Comissão Eleitoral em exercício.
}
\todo[inline]{\textbf{Comentário:} Considerar se "2 (dois) semestres letivos inteiros sem Diretoria" pode ser interpretado de forma maliciosa por insatisfação com gestões "pouco ativas".}
\paragrafo{As eleições emergenciais decorrentes de vacância total deverão ser obrigatoriamente concluídas no prazo máximo de 45 (quarenta e cinco) dias corridos a partir da formação da Comissão Eleitoral na Assembleia.}
\paragrafo{A contagem deste prazo será suspensa durante os períodos de recesso e férias acadêmicas da UFMG, bem como durante os períodos de restrição de início e fim do semestre letivo, conforme disposto no Art. [46] deste Estatuto.}
\paragrafo{Em caso de greve de servidores ou professores da UFMG durante este prazo, a Assembleia Geral poderá, se entender necessário, suspender a contagem enquanto a greve durar.}
\todo[inline,backgroundcolor=blue!70]{\textbf{Referência:} Art. [46] – O processo eleitoral deverá ocorrer integralmente dentro do semestre letivo [...] O processo deve iniciar no mínimo 30 (trinta) dias após o início [...] e encerrar com antecedência mínima de 30 (trinta) dias do término.}
\paragrafo{Os membros desta Comissão Eleitoral excepcional ficam expressamente proibidos de concorrer ao pleito que conduzirem.}
}

%===================== SEÇÃO II =====================

\secao{Da Assembleia Geral}

\artigo{
A Assembleia Geral é composta por todos os associados em pleno gozo de seus direitos.

\paragrafoUnico{
A Assembleia Geral tem como finalidade debater assuntos de interesse do Corpo Discente e elaborar propostas a serem submetidas a plebiscito, não possuindo caráter deliberativo, salvo nas exceções expressamente previstas neste Estatuto.
\todo[inline,backgroundcolor=purple!30]{\textbf{Update:} A ressalva para exceções expressamente previstas no Estatuto foi adicionada para casos como a Vacância Total.}
\todo[inline,backgroundcolor=blue!70]{\textbf{Referência:} "Em caso de vacância total da Diretoria Executiva, uma Assembleia Geral Extraordinária deverá ser convocada para, em caráter deliberativo excepcional e com dispensa de plebiscito, formar uma Comissão Eleitoral (...)" / "Deliberar, em caráter excepcional e com dispensa de plebiscito, sobre a formação de Comissão Eleitoral em casos de vacância total da gestão, conforme os termos previstos neste Estatuto."}
\todo[inline]{\textbf{Comentário:} O caráter não deliberativo é importante para evitar táticas maliciosas (por exemplo, marcar em dia inconveniente, impor voto presencial, ou escolher horários para a votação que favoreçam um grupo específico). Assim, permanece apenas o debate.}
}
}

\todo[inline,backgroundcolor=purple!30]{\textbf{Update:} Definições dos tipos de Assembleia movidas para cá (antes ficavam após as regras de prazos e quóruns).}
\artigo{
As Assembleias podem ser:
\inciso{
\item Ordinárias;
\item Extraordinárias.
}
}

\artigo{
A Assembleia Ordinária ocorrerá anualmente, em até cinco dias úteis após as eleições da Diretoria Executiva, com a finalidade exclusiva de:
\inciso{
\item Discutir e esclarecer dúvidas sobre a prestação de contas financeira e patrimonial da Diretoria anterior;
\item Avaliar o relatório da Comissão Eleitoral, podendo, caso constate irregularidades ou fraudes eleitorais graves, decidir pela retenção cautelar da posse e encaminhar proposta de anulação de eleição para o Plebiscito;
\item Dar posse aos membros eleitos, caso o pleito seja validado.
}
}
  
\artigo{
A Assembleia Extraordinária somente poderá ser convocada pela Diretoria Executiva ou por requerimento de, no mínimo, 10\% (dez por cento) dos associados, respeitado o caráter não deliberativo e os prazos regimentais de realização.
\todo[inline]{\textbf{Comentário:} O prazo máximo de 15 dias, agora centralizado nas regras de Assembleia, evita a inação das partes.}
\todo[inline,backgroundcolor=blue!70]{\textbf{Referência:} - Parágrafo Único - A Assembleia Geral tem como finalidade debater [...] não possuindo caráter deliberativo.}
}

\artigo{
A convocação e a realização das Assembleias Gerais padronizam-se pelas seguintes regras de prazos, salvo disposição em contrário:
\inciso{
\item Todo edital de convocação deverá ser publicado com antecedência mínima de 6 (seis) dias úteis da data da Assembleia, por meio dos canais digitais oficiais do \textit{\SiglaDoCA{}} e em locais de ampla circulação;
\item Em casos excepcionais de urgência, fundamentados exclusivamente na ocorrência de fato novo e relevante detalhado no próprio edital e obrigatoriamente enquadrado no parágrafo único deste artigo, a antecedência do edital poderá ser reduzida para até 72 (setenta e duas) horas;
\item Quando a Assembleia Extraordinária for solicitada por meio de requerimento discente ou em virtude da tramitação de recurso disciplinar, a sua realização não poderá ultrapassar o prazo máximo de 15 (quinze) dias corridos, contados a partir da data de entrega formal do pedido à Diretoria Executiva.
}
\paragrafoUnico{Para fins de redução excepcional de prazo, compreendem-se como fato novo e relevante tão somente as situações que decorram de decisão administrativa ou acadêmica da Universidade com impacto imediato sobre os associados, que envolvam risco atestado de perda de prazo institucional ou que representem ameaça concreta aos direitos estabelecidos neste Estatuto.}
\todo[inline]{\textbf{Comentário:} É um trecho sensível, pois pode ser usado para justificar convocação de assembleias em momentos inconvenientes. Restringir à ameaça a direitos do estatuto restringe bem o escopo. Decisão da universidade também parece suficiente. Apenas a não especificação do que é "perda de prazo institucional" que pode ter uma abrangência imprevista.}
}

\artigo{
A Assembleia será instalada:
\inciso{
\item Em primeira convocação, com a presença mínima de 5\% (cinco por cento) dos associados;
\item Em segunda convocação, 20 (vinte) minutos após, com qualquer quórum presente;
\item Para assuntos sensíveis, compreendendo alteração do Estatuto, destituição total ou parcial de membros da Diretoria Executiva, anulação de processo eleitoral e dissolução do \textit{\SiglaDoCA{}}, será exigida a presença mínima de 10\% (dez por cento) dos associados ou de 50 (cinquenta) alunos, o que for menor, em quórum rígido e sem segunda convocação.
}
}

\artigo{
Todo edital de convocação de Assembleia deverá indicar de forma clara e específica:
\inciso{
\item a pauta detalhada dos assuntos a serem discutidos;
\item a data, o horário e, quando aplicável, o local físico da reunião;
\item os meios de participação e autenticação.
}
}

\artigo{
A Assembleia somente poderá discutir os assuntos e votar os respectivos encaminhamentos e propostas para plebiscito expressamente indicados na pauta do edital de convocação, sendo nulos quaisquer deliberações sobre temas não previstos.
\todo[inline]{\textbf{Comentário:} Buscando evitar a possível manipulação de propostas ao colocar pauta genérica. A assembleia geral que deflagrou a "greve dos estudantes" em 2024 não tinha nenhuma declaração nos posts do Instagram dizendo que era para votar uma greve estudantil. O caráter não deliberativo da assembléia também diminui este risco.}
}

\artigo{
Compete à Assembleia:
\inciso{
\item Discutir e elaborar propostas a serem submetidas a plebiscito;
\item Elaborar propostas de reforma do Estatuto, destituição de membros da Diretoria Executiva, anulação de processo eleitoral ou dissolução do \textit{\SiglaDoCA{}}, a serem submetidas a plebiscito exclusivo;
\item Emitir pareceres sobre casos omissos neste Estatuto, a serem submetidos a plebiscito;
\item Discutir e elaborar pareceres sobre recursos relativos a penalidades disciplinares ou a decisões da Diretoria ou da própria Assembleia que tenham violado direitos de associados, a serem submetidos a Plebiscito;
\item Deliberar, em caráter excepcional e com dispensa de plebiscito, sobre a formação de Comissão Eleitoral em casos de vacância total da gestão, conforme os termos previstos neste Estatuto.
\todo[inline,backgroundcolor=purple!30]{\textbf{Update:} Esta competência deliberativa foi adicionada exclusivamente para o caso de vacância total, com a ressalva de seguir os termos previstos no Estatuto para reforçar o amparo legal da exceção.}
\todo[inline,backgroundcolor=blue!70]{\textbf{Referência:} "Em caso de vacância total da Diretoria Executiva, uma Assembleia Geral Extraordinária deverá ser convocada para, em caráter deliberativo excepcional e com dispensa de plebiscito, formar uma Comissão Eleitoral (...)"}
}
}

\artigo{
Os pareceres e propostas formulados pela Assembleia, acompanhados de suas respectivas justificativas, bem como da lista nominal dos presentes,
deverão ser divulgados em até 3 (três) dias úteis após sua realização, por meio dos canais digitais oficiais do \textit{\SiglaDoCA{}}.
\paragrafoUnico{As contribuições em Assembleia deverão ser registradas em ata pública.}
}

%===================== SEÇÃO III =====================
\secao{Do Plebiscito}

\artigo{
O plebiscito é o instrumento deliberativo e de decisão final do \textit{\SiglaDoCA{}}, destinado a aprovar ou rejeitar propostas de interesse do Corpo Discente.
\paragrafoUnico{As propostas elaboradas pela Assembleia somente terão efeito após aprovação em Plebiscito.}
}

\artigo{
O plebiscito poderá ser convocado e organizado:
\inciso{
\item pela Diretoria Executiva, por iniciativa própria ou de Assembleia por ela convocada;
\item por Comissão Organizadora designada em Assembleia convocada por requerimento discente, nos termos do Art. 20.
}
\todo[inline,backgroundcolor=blue!70]{\textbf{Referência:} Art. Xº – A Assembleia Extraordinária somente poderá ser convocada pela Diretoria [...]}
}

\artigo{
A Comissão Organizadora será composta por, no mínimo, 2 (dois) associados em pleno gozo de seus direitos, escolhidos no momento da convocação ou da Assembleia, com mandato restrito à condução do plebiscito específico.
}

\artigo{
O edital de convocação do plebiscito deverá conter:
\inciso{
\item a redação exata da proposta ou pergunta submetida;
\item a justificativa correspondente;
\item as datas, os horários e a duração da votação;
\item os locais e horários para votação presencial, quando houver;
\item os meios de participação e autenticação.
}
}

\artigo{
A duração do plebiscito será de, no mínimo, 2 (dois) dias úteis, garantindo-se o funcionamento ininterrupto da votação em meios remotos ou digitais e, se houver modalidade presencial, o funcionamento em horários que abranjam os turnos letivos dos associados.
\todo[inline,backgroundcolor=purple!30]{\textbf{Update:} Artigo desmembrado para separar as regras de conteúdo do edital das regras de duração do plebiscito.}
}

\artigo{
O edital deverá ser publicado com antecedência mínima de 7 (sete) dias úteis.
\paragrafoUnico{Em casos de urgência justificada,
\todo[inline]{\textbf{Comentário:} Pensar em como funcionariam os prazos para o caso de tentar interromper eleições.}
o prazo poderá ser reduzido para até 72 (setenta e duas) horas, nos termos das regras de convocação de Assembleia dispostas no Art. [21].
\todo[inline,backgroundcolor=blue!70]{\textbf{Referência:} Art. [21] - A convocação e a realização das Assembleias Gerais padronizam-se [...] Em casos excepcionais de urgência, fundamentados exclusivamente na ocorrência de fato novo e relevante [...] a antecedência do edital poderá ser reduzida para até 72 (setenta e duas) horas.}
}
}

\artigo{
O plebiscito será válido com participação mínima de 10\% (dez por cento) dos associados ou 100 (cem) votantes, o que for menor.
\todo[inline]{\textbf{Comentário:} Em regra, não seria necessário quórum.}
\paragrafo{Para reforma do Estatuto, destituição de membros da Diretoria Executiva, anulação de processo eleitoral ou dissolução do Centro Acadêmico, será exigida a participação mínima de 25\% (vinte e cinco por cento) dos associados.}
\paragrafo{As decisões ordinárias serão aprovadas por maioria simples dos votos válidos.}
\paragrafo{As matérias críticas dependerão de aprovação por maioria qualificada de 2/3 (dois terços) dos votos válidos.
\todo[inline]{\textbf{Comentário:} Especificar o que será considerado matéria crítica.}}
\paragrafo{Antes de declarar a invalidade por falta de quórum em matérias críticas, a votação poderá ser prorrogada por até 48 (quarenta e oito) horas, mediante nova comunicação pública.}
}

\artigo{
O voto no plebiscito será direto e facultativo, podendo ser secreto ou não, conforme definido no edital, de acordo com a natureza da matéria.
}

\artigo{
A apuração será imediata após o encerramento da votação e os resultados publicados em até 3 (três) dias úteis, indicando o total de votantes, votos válidos, brancos e nulos, acompanhados da ata digital acessível a todos os associados.
}

\artigo{
Em caso de falha técnica grave ou fraude comprovada, a instância responsável pela condução do plebiscito, seja a Diretoria Executiva ou a Comissão Organizadora, poderá suspender e remarcar a votação, mediante decisão fundamentada.
\paragrafo{A remarcação deverá ocorrer em até 10 (dez) dias úteis.}
\paragrafo{Recursos poderão ser apresentados por qualquer associado em até 3 (três) dias úteis após a divulgação dos resultados, devendo a instância responsável (Diretoria ou Comissão) decidir em até 5 (cinco) dias úteis.}
\paragrafo{É vedada a interferência da Diretoria Executiva em plebiscitos conduzidos por Comissão Organizadora, salvo para garantir suporte técnico ou logístico solicitado pela instância responsável.
\todo[inline]{\textbf{Comentário:} Exagero?}}
}


%===================== CAPÍTULO IV =====================
\capitulo{DO PROCESSO ELEITORAL}

\artigo{
As eleições para a Diretoria Executiva ocorrerão anualmente, no segundo semestre letivo da UFMG, seguindo o processo eleitoral deste capítulo.
\todo[inline]{\textbf{Comentário:} Em universidades públicas o calendário civil (janeiro a dezembro) frequentemente descola do calendário letivo devido a paralisações.}
\paragrafoUnico{A Diretoria Executiva poderá antecipar as eleições e o término do seu mandato.}
}

\todo[inline,backgroundcolor=purple!30]{\textbf{Update:} Artigo sobre o formato do voto e quórum eleitoral movido para o início do Capítulo IV (antes localizado no final do capítulo), a fim de estabelecer as premissas da eleição logo em sua abertura.}
\artigo{
As eleições serão realizadas pelo voto direto, secreto e facultativo.
\paragrafoUnico{Será considerada inválida a eleição com participação de menos de 5\% (cinco por cento) dos associados ou de menos de 40 participantes, o que for maior.
\todo[inline]{\textbf{Comentário:} Atualmente (2025.2): (EFTCC: 197) (BCC: 400) (BSI: 386) [15\% de BCC: 60] [10\% de BCC + EFTCC: 60]}
\todo[inline]{\textbf{Comentário:} Outras universidades, como a UNIFEI, possuem a regra de que é necessária a participação de 1/3 dos alunos para eleger diretoria. E funciona, pois são obrigados a pedir bastante aos alunos, e caso não consigam, a eleição deve se repetir.}
}
}

%\artigo{ As eleições para a Representação Discente em Órgãos Colegiados da UFMG ocorrerão sempre quando os respectivos órgãos solicitarem ou quando houver vacância de cargo, seguindo o processo eleitoral deste capítulo. 
%\paragrafoUnico{A eleição para Representação Discente poderá ser dispensada apenas quando a pessoa representante titular abdicar ou perder o seu mandato, permitindo uma indicação temporária e emergencial por, no máximo, 45 (quarenta e cinco) dias. \todo[inline]{Bomba nuclear, precisa pensar. Precisa exigir que a explicação seja publicada. Está em lugar péssimo, mas criar seção para representação (assim como diretoria) vale o espaço? “Perder o mandato” parece perigoso. Está claro?}}}

\artigo{
As eleições para a Representação Discente em Órgãos Colegiados da UFMG ocorrerão sempre que houver vacância do cargo, seguindo o processo eleitoral deste Estatuto.
\todo[inline]{\textbf{Comentário:} Avaliar ambiguidade do "perder o mandato".}
\todo[inline]{\textbf{Comentário:} Incluir, nos direitos, o direito de conhecer quem são as pessoas representantes nos órgãos colegiados.}
\paragrafo{A eleição poderá ser dispensada exclusivamente em caso de renúncia ou perda de mandato, permitindo à Diretoria Executiva realizar indicação temporária por até 45 (quarenta e cinco) dias corridos, prazo improrrogável no qual novas eleições devem ser obrigatoriamente concluídas.}
\paragrafo{Caso não haja candidaturas no processo eleitoral, a Diretoria Executiva poderá realizar indicação direta, válida até o fim do seu mandato.}
}

\todo[inline,backgroundcolor=purple!30]{\textbf{Update:} Artigo sobre o mandato da representação discente movido para logo em seguida da regra sobre o seu respectivo rito de eleição.}
\artigo{
O mandato do representante discente se dará de acordo com as normas da UFMG e do respectivo Órgão Colegiado.
}

\artigo{
A Comissão Eleitoral será composta por 1 (um) presidente e 1 (um) secretário, ambos indicados pela Diretoria Executiva.
\paragrafo{É vedada a composição da Comissão Eleitoral por membros da atual Diretoria Executiva ou integrantes de chapas inscritas.}
\paragrafo{É assegurado a cada chapa inscrita o direito de indicar 1 (um) representante para fiscalizar os trabalhos da Comissão, sem poder de voto.}
}

\artigo{
Será função da Comissão Eleitoral garantir a lisura e isonomia do processo eleitoral, responsabilizando‑se por todas as etapas do processo, sempre observando as disposições deste estatuto.
\paragrafo{A Comissão Eleitoral possui caráter executório, resolvendo as questões do dia a dia da eleição.}
\todo[inline]{\textbf{Comentário:} Trava de delimitação: a Comissão não legisla, ela apenas executa o edital.}
\paragrafo{Fica proibido à Comissão Eleitoral criar ou alterar regras com o pleito em andamento. Em especial, é vedado inovar em métodos de apuração, critérios para validar ou anular votos, e demais normas que definam os vencedores, caso não estejam expressamente descritos neste Estatuto.}
\todo[inline]{\textbf{Comentário:} Garante estabilidade ao fechar a brecha que permitia à Comissão inventar regras para manipular a eleição no último momento.}
}

\artigo{
A Comissão Eleitoral é responsável única e exclusivamente pelo processo eleitoral, não tendo nenhum poder ou condição de equivalência com a Diretoria Executiva.
\paragrafoUnico{Caso aconteça algum problema grave, fraude ou situação não prevista por este Estatuto, a Comissão Eleitoral não poderá suspender ou anular a eleição por conta própria. Ela deve registrar tudo em ata, garantir que o processo ocorra normalmente até o fim, terminar a apuração, e encaminhar o problema para avaliação da Assembleia Ordinária subsequente.}
\todo[inline]{\textbf{Comentário:} Impede a "política do fato consumado". Evita que a Comissão interrompa o processo eleitoral sem transparência, transferindo todo caso atípico para ser julgado publicamente na Assembleia.}
}

\artigo{
Durante o processo eleitoral para a Diretoria Executiva, as funções da Diretoria Executiva ficam suspensas.
\paragrafo{A Diretoria Executiva deverá entregar à Comissão Eleitoral todas as chaves, credenciais, recursos financeiros, meios de comunicação e demais instrumentos de gestão sob sua guarda ou controle.
\todo[inline]{\textbf{Comentário:} Evita favorecimento durante o processo eleitoral. Mas será que há algum risco de ninguém poder fazer nada? Há algo relevante além da indicação de pessoas representantes em órgãos colegiados?}
}
\paragrafo{A Diretoria Executiva somente retomará suas funções após a posse da nova Diretoria eleita.}
\paragrafo{Excepcionalmente, a Diretoria Executiva poderá realizar os trâmites necessários para garantir a representação discente nos órgãos colegiados, mediante aprovação prévia da Comissão Eleitoral.}
}

\artigo{
O processo eleitoral será iniciado pela Diretoria Executiva por meio de edital a ser amplamente divulgado com antecedência mínima de 15 (quinze) dias antes da data de início das eleições.
\paragrafoUnico{Deverá ser respeitado o prazo de 7 (sete) dias para a inscrição das chapas que concorrerão nas eleições.}
}

\artigo{
O processo eleitoral deverá ocorrer integralmente dentro do semestre letivo, respeitando os seguintes prazos:
\inciso{
\item O processo deve iniciar no mínimo 30 (trinta) dias após o início do semestre letivo da UFMG;
\item O processo deverá encerrar com antecedência mínima de 30 (trinta) dias do término do semestre letivo da UFMG.
}
}

\artigo{
A eleição se dará por chapas, sendo vencedora a chapa que obtiver o maior número de votos, não computados os nulos e os em branco.
\paragrafo{Qualquer associado poderá integrar uma chapa, desde que não esteja em situação de trancamento total do curso, durante o semestre em que ocorrerem as eleições.}
\paragrafo{A composição das chapas, assim como a documentação delas, não poderá ser alterada ou substituída após o término do período de inscrição.}
\todo[inline,backgroundcolor=purple!30]{\textbf{Update:} Antes estava "alterada, suprimida ou substituída".}
\paragrafo{Para as eleições da Representação Discente em órgãos colegiados, as chapas serão consideradas por cadeira (um titular e um suplente).}
\paragrafo{Para as eleições da Representação Discente, não será permitida a inscrição de chapas que incluam associados que já ocupem cargo de Representação Discente ou que tenham se candidatado a outra vaga de representação.}
\todo[inline]{\textbf{Comentário:} Como evitar que uma Comissão Eleitoral decida, sem respaldo, fazer com que se o número de votos nulos ultrapassar o número de votos para uma outra chapa (mesmo que seja em chapa única)? Foi algo que já aconteceu.}
}

\artigo{
Fica impugnada a chapa para a Diretoria Executiva que possuir, em sua composição, um ou mais membros de qualquer diretoria (atual ou passadas) que ainda não tenha publicado a sua respectiva prestação de contas financeira até o término do período de inscrição das chapas, conforme exigido no presente Estatuto.
\todo[inline]{\textbf{Comentário:} A regra impede que estudantes com histórico de omissão com as finanças de suas respectivas gestões voltem a se eleger até que regularizem as contas do seu período.}
\todo[inline]{\textbf{Comentário:} Não existem meios simples de se punir uma diretoria por se recusar a fazer alguma coisa que ela tem o dever de fazer. Este é um exemplo, mas ele é apenas relevante para alguém que ainda queira continuar se candidatando.}
}

%\artigo{A eleição será anulada se a diferença entre o número de votantes registrados e o de votos ultrapassar 5\% (cinco por cento) do total de votantes, ou se for suficiente para alterar o resultado do pleito.\paragrafoUnico{Em hipótese nenhuma a eleição poderá ser interrompida, exceto por decisão de plebiscito específico para esse propósito.}}

\artigo{
A eleição será automaticamente anulada pela Comissão Eleitoral caso a diferença absoluta entre o número de votantes registrados e o de votos apurados se enquadre em ao menos uma das seguintes situações:
\todo[inline]{\textbf{Comentário:} A Comissão deveria ter autonomia para anular por outro fundamento?}
\todo[inline]{\textbf{Comentário:} Mecanismo anti-fraude, caso seja escolhido oferecer a opção do voto presencial (urna).}
\inciso{
\item for superior a 5\% (cinco por cento) do total de votantes registrados; ou
\item for igual ou superior à diferença de votos entre a chapa primeira colocada e a chapa segunda colocada.
}
}


\artigo{
O processo eleitoral não poderá ser interrompido ou suspenso em nenhuma hipótese, inclusive em decorrência da convocação prévia ou intercorrente de Assembleias ou Plebiscitos.
\paragrafoUnico{Quaisquer suspeitas de fraudes ou irregularidades que não gerem a anulação automática prevista neste Estatuto deverão constar em relatório da Comissão Eleitoral e ser julgadas exclusivamente na Assembleia Ordinária de posse.}
\todo[inline,backgroundcolor=purple!30]{\textbf{Update:} Regra de ininterrupção do pleito separada do artigo sobre anulação matemática por fraude, garantindo foco único por artigo.}
\todo[inline]{\textbf{Comentário:} Não se admite interrupção do pleito. Qualquer anulação interpretativa (não matemática) deve ser debatida na Assembleia Ordinária após o encerramento da votação.}
}

\artigo{
Em caso de anulação definitiva do processo eleitoral aprovada em Plebiscito, a cerimônia de posse prevista não terá efeito.
\paragrafoUnico{A mesma Assembleia Geral Extraordinária que encaminhou a proposta de anulação designará a nova Comissão Eleitoral, cujas funções, restrições e prazos seguirão as mesmas regras aplicáveis aos casos de vacância total da Diretoria Executiva.
\todo[inline,backgroundcolor=blue!70]{\textbf{Referência:} Regras de Vacância Total - As eleições deverão ser obrigatoriamente concluídas no prazo máximo de 45 dias... a contagem será suspensa nas férias e recessos... os membros ficam proibidos de concorrer.}
}
\todo[inline]{\textbf{Comentário:} Remove a subjetividade na escolha dos membros e foca o grupo apenas em realizar novas eleições, garantindo que os envolvidos nas chapas ou comissões do pleito anulado sofram as restrições da vacância.}
}

%\artigo{A eleição será anulada se a diferença absoluta entre o número de votantes registrados e o total de votos apurados ultrapassar 5\% (cinco por cento) dos votantes, ou se essa diferença for matematicamente suficiente para alterar o resultado do pleito.\paragrafoUnico{Em hipótese nenhuma a eleição poderá ser interrompida, exceto por decisão de plebiscito específico para esse propósito.}}

\artigo{
A apuração dos votos deverá ser feita pela Comissão Eleitoral imediatamente após o término da votação, a fim de garantir a exatidão dos resultados.
}

%\artigo{O mandato da Diretoria Executiva será de no máximo 12 (doze) meses, ou até o próximo segundo semestre do ano letivo da UFMG, o que for mais curto.\todo[inline]{Em universidades públicas o calendário civil (janeiro a dezembro) frequentemente descola do calendário letivo devido a paralisações.}\paragrafoUnico{Se a UFMG estiver em greve de funcionários e/ou professores, durante o período de eleições, o mandato da Diretoria Executiva poderá ser prorrogado até o fim da greve.}\todo[inline]{Fica a cargo da Diretoria Executiva decidir a prorrogação ou não. Em greves de professores, é inviável realizar eleições. Mas em alguns casos, greves de servidores podem não interromper a rotina da maior parte dos alunos e o funcionamento da universidade.}}

\artigo{
Se a UFMG estiver em greve de funcionários e/ou professores, durante o período previsto para as eleições, o mandato da Diretoria Executiva poderá ser prorrogado até 20 (vinte) dias corridos após o fim da greve.
\todo[inline]{\textbf{Comentário:} Fica a cargo da Diretoria Executiva decidir a prorrogação ou não. Em greves de professores, é inviável realizar eleições. Mas em alguns casos, greves de servidores podem não interromper a rotina da maior parte dos alunos e o funcionamento da universidade.}
}


%===================== CAPÍTULO V =====================
\capitulo{DO PATRIMÔNIO E DAS FINANÇAS}

\artigo{
O patrimônio do \textit{\SiglaDoCA{}} será constituído por bens móveis, imóveis, doações, legados e recursos financeiros provenientes de:
\inciso{
\item Contribuições voluntárias dos associados;
\item Renda de eventos e campanhas promovidas pelo \textit{\SiglaDoCA{}};
\item Doações e patrocínios de terceiros.
}
}

\artigo{
Os recursos financeiros serão depositados e movimentados em conta bancária, em nome do \textit{\SiglaDoCA{}}, sob a responsabilidade conjunta da Presidência e da Tesouraria.
\paragrafoUnico{A exigência de que a conta bancária esteja em nome do \textit{\SiglaDoCA{}} poderá ser dispensada, caso se comprove a impossibilidade de sua abertura em nome da entidade por questões legais ou administrativas.}
}

\artigo{
A prestação de contas financeira da gestão é obrigatória e deve ser materializada por meio de um relatório financeiro que detalhe as receitas, as despesas e o saldo em caixa, de fácil compreensão aos associados.
\paragrafo{A elaboração do relatório e a sua respectiva publicação nos canais de comunicação do \textit{\SiglaDoCA{}} são responsabilidades da Tesouraria ou, na omissão desta, da Presidência.}
\paragrafo{A publicação do relatório deverá ocorrer, impreterivelmente, até o término do período de inscrição de chapas do processo eleitoral.}
\paragrafo{O relatório publicado deverá ser submetido à Assembleia Geral Ordinária.}
}

\artigo{
Em caso de dissolução do \textit{\SiglaDoCA{}}, o seu patrimônio remanescente será destinado a uma entidade congênere, a ser proposta pela Assembleia Geral e aprovada em Plebiscito convocado para deliberar sobre a dissolução.
}

%===================== CAPÍTULO VI =====================
\capitulo{DAS DISPOSIÇÕES GERAIS E TRANSITÓRIAS}

\artigo{
O presente Estatuto só poderá ser alterado mediante aprovação em Plebiscito, a partir de proposta formulada em Assembleia Geral Extraordinária específica para este fim.
}

\artigo{
Os membros da Diretoria Executiva não serão remunerados pelo exercício de suas funções.
}

\artigo{
Os casos omissos neste Estatuto serão resolvidos pela Diretoria Executiva.
\paragrafo{Qualquer decisão sobre caso omisso deverá ser amplamente divulgada, cabendo recurso à Assembleia Geral, que poderá formular proposta de veto a ser deliberada em Plebiscito.}
\paragrafo{As decisões sobre casos omissos terão validade restrita ao mandato da Diretoria Executiva que as editou, salvo se posteriormente incorporadas a este Estatuto mediante reforma aprovada em Plebiscito.}
}

\artigo{
Das decisões de cada diretoria cabe recurso com efeito resolutivo à reunião da Diretoria Executiva. Das decisões da Diretoria Executiva cabe recurso à Assembleia Geral, que elaborará parecer ou proposta a ser submetida a plebiscito, instância final de deliberação do Corpo Discente.
}

\artigo{
Cabe à Diretoria Executiva zelar pelo cumprimento do presente Estatuto.
}

\todo[inline,backgroundcolor=purple!30]{\textbf{Update:} Adicionado para permitir a fundação direta por plebiscito assíncrono, contornando a dificuldade de quórum de uma assembleia.}
\artigo{
Excepcionalmente, a fundação do \textit{\SiglaDoCA{}}, a aprovação fundamental deste Estatuto e o desmembramento da representação estudantil antecedente dar-se-ão de forma direta por meio de Plebiscito de Fundação, dispensando-se a exigência de Assembleia Geral prévia.
\paragrafo{A legitimidade da convocação do Plebiscito de Fundação advém da presidência da representação estudantil unificada antecedente.}
\paragrafo{A substituição do debate síncrono em Assembleia dar-se-á por um período prévio de consulta e discussão pública assíncrona, no qual o projeto de estatuto será disponibilizado aos estudantes afetados para o envio de propostas e emendas antes do fechamento do texto para votação.}
\paragrafo{A Comissão Organizadora do Plebiscito de Fundação será designada pela autoridade convocante, sendo responsável por definir e dar publicidade aos prazos e meios de votação.}
}

\artigo{
Este Estatuto entrará em vigor na data de sua aprovação em Plebiscito, revogadas quaisquer disposições anteriores que conflitem com este.
}

%===================== ASSINATURAS =====================
% Assinaturas (quando o Estatuto for aprovado em Plebiscito)
% Assinaturas das pessoas responsáveis pela condução do plebiscito


\begin{center}
\vspace{1.5cm}
\noindent\textbf{Belo Horizonte, Abril de 2026}
\vspace{1.5cm}

\begin{tabular}{c c}
\fbox{\vphantom{\rule{0pt}{2.2cm}}\hspace{6.0cm}} &
\fbox{\vphantom{\rule{0pt}{2.2cm}}\hspace{6.0cm}} \\
\\
\begin{tabular}{c}
Nome da Pessoa \\
\small Presidente da Comissão Organizadora
\end{tabular}
&
\begin{tabular}{c}
Nome da Pessoa \\
\small Secretário(a) da Comissão Organizadora
\end{tabular}
\end{tabular}

\vspace{1.2cm}

% linha de baixo com caixa do mesmo tamanho, centralizada usando uma tabela de três colunas
\begin{tabular}{c c c}
 & \fbox{\vphantom{\rule{0pt}{2.2cm}}\hspace{6.0cm}} &  \\
 & &  \\
 & 
 \begin{tabular}{c}
Nome da Pessoa \\
 \small Membro da Comissão Organizadora
 \end{tabular}
 & 
\end{tabular}
\end{center}

\end{document}